
On a alors, d'après IV (2.12.4) :
@@M0@@
la dernière égalité résultant de ce que pour @@M1@@, on a @@M2@@ (VI 2.8). Supposons que @@M3@@. Alors ce qui précède montre que @@M4@@; appliquant le même raisonnement aux éléments @@M5@@ pour @@M6@@, on en déduit par récurrence que les @@M7@@ sont tous nuls, d'où l'injectivité.

3.3. Soient toujours @@M8@@ un schéma noethérien, et @@M9@@ un @@M10@@-Module localement libre de rang @@M11@@. On note @@M12@@ le fibré vectoriel (contravariant en @@M13@@!) associé à @@M14@@ (SGA 3 I 4.6), considéré comme schéma en groupes sur @@M15@@. Soit @@M16@@ un torseur sous @@M17@@; on se propose de calculer @@M18@@.

Rappelons comment on peut plonger @@M19@@ dans un fibré projectif sur @@M20@@. Les torseurs sous @@M21@@ correspondent aux extensions @@M22@@ de @@M23@@ par @@M24@@. En effet, si on a une extension :
@@M25@@
on en déduit un torseur sous @@M26@@ en considérant l'homomorphisme @@M27@@ défini par @@M28@@, la section de @@M29@@ définie par @@M30@@, et en prenant @@M31@@, les opérations de @@M32@@ sur @@M33@@ étant induites par celles de @@M34@@ sur @@M35@@.

En sens inverse, soit @@M36@@ un torseur sous @@M37@@. On note @@M38@@ le faisceau des germes d'homomorphismes de @@M39@@ dans @@M40@@; c'est un faisceau principal homogène sous le faisceau des germes d'homomorphismes de @@M41@@ dans @@M42@@, qui s'identifie au dual @@M43@@ de @@M44@@. On en déduit alors de façon classique une extension de @@M45@@ par @@M46@@, d'où en passant aux duaux la suite (3.3.1).

